\documentclass[addpoints]{exam}
% \documentclass[addpoints,12pt]{exam}

\usepackage[slovene]{babel}
\usepackage{amsfonts}
\usepackage{amssymb}
\usepackage{amsmath}
\usepackage{float}
\usepackage{graphicx}
\usepackage{enumitem}
\usepackage{color}
\usepackage[gen]{eurosym}
\usepackage{newunicodechar}
\newunicodechar{€}{\euro}

%Komentar naslednje vrstice glede na verzijo testa -- rešitve ali prostor za odgovore:
% \printanswers

% \linespread{2}


\pointsinrightmargin
\marginbonuspointname{}


\renewcommand{\solutiontitle}{\noindent\textbf{Rešitve in točkovanje:}\par\noindent}

\colorgrids
\definecolor{GridColor}{gray}{0.7}
\setlength{\gridsize}{7mm}

\extrawidth{5mm}
\setlength{\rightpointsmargin}{1.5cm}

\begin{document}

\runningfooter{}{}{}

\begin{center}
    \fbox{\fbox{\parbox{15cm}{\centering
    Popravljanje in pridobivanje ocen -- konferenca \\ 1. letnik -- test C \\ 15. junij 2026 \\ (Realna števila; Koordinatni sistem v ravnini; Funkcija; Premica)}}}
\end{center}

\vspace{0.1in}
\makebox[\textwidth]{Ime in priimek, oddelek:\enspace\hrulefill}


\begin{center}
    \hqword{Naloga}
    \hpword{Možne točke}
    \chbpword{Dodatne točke}
    \hsword{Dosežene točke}
    \htword{Skupaj}  
    % \combinedpointtable[h][questions]
    \gradetable[h][questions]

\end{center}

\vspace{0.02in}


\begin{questions}

%%%% realna števila


    \question[5]
    Izračunajte natančno vrednost izraza.
    $$ \sqrt{6\cdot 8\cdot 12}-\left(\sqrt{3}-1\right)^2+\dfrac{\sqrt{3}}{2+\sqrt{3}}-\left(3+\sqrt{3}\right)\left(1-2\sqrt{3}\right) $$

    \begin{solution}[8cm]
        \begin{align*} 
            \dfrac{\sqrt{5}}{\sqrt{5}+1}-\dfrac{5-3\sqrt{5}}{4\sqrt{5}}+\left(3+\sqrt{5}\right)^2-\sqrt{125}&=\dfrac{\sqrt{5}(\sqrt{5}-1)}{(\sqrt{5}+1)(\sqrt{5}-1)}-\dfrac{(5-3\sqrt{5})\sqrt{5}}{4\sqrt{5}\sqrt{5}}+9+6\sqrt{5}+5-5\sqrt{5}= \\
                                                                                                            &=\dfrac{5-\sqrt{5}}{5-1}-\dfrac{5\sqrt{5}-3\cdot 5}{4\cdot 5}+14+\sqrt{5}= \\
                                                                                                            &=\frac{5-\sqrt{5}}{4}-\dfrac{\sqrt{5}-3}{4}+14+\sqrt{5}= \\
                                                                                                            &=\frac{8-2\sqrt{5}}{4}+14+\sqrt{5}= \\
                                                                                                            &=16+\frac{1}{2}\sqrt{5}
        \end{align*}
        Za racionalizacijo $2$ točki, za izračun kvadrata dvočlenika $1$ točka, delno korenjenje $1$ točka, izračun $1$ točka, končne rezultat $1$ točka.
    \end{solution}
    
    

    % 7. vprašanje
    \question[5]
    Rešite enačbo in opišite njen geometrijski pomen ter rešitev grafično prikažite.
    $$  \left\lvert x+2\right\rvert=8 $$
    
    
    \begin{solution}[5cm]
        Opis geometrijskega pomena (oddaljenost števila $x$ od števila $-3$ je $2$) $1$ točka, grafična predstavitev $1$ točka, reševanje enačbe z upoštevanjem pogojev do $2$ točke, zapis rešitve $x\in\{-5,-1\}$ $1$ točka..
    \end{solution}
        



    



    \newpage 

    
    % 3. vprašanje
    \question[5]
    Rešite sistem neenačb, rešitev grafično upodobite in zapišite z intervalom.
    $$ \left(3x-5>\frac{2}{3}x+2\right)\land\left(\frac{2}{5}x\geq x-3\right) $$
    
    
    \begin{solution}[10cm]
        Rešitev prve enačbe $x>2$ ali $x\in(2,\infty)$ $1$ točka, rešitev druge enačbe $x\geq 3$ ali $x\in[3,\infty)$ $1$ točka, rešitev sistema $x\geq 3$ ali $x\in[3,\infty)$ $1$ točka, zapis končne rešitve z intervalom $1$ točka, grafična upodobitev $1$ točka.
    \end{solution}





    % 9. vprašanje
    \question[5]
    Prvi delavec sam pozida steno v $10$ urah, drugi v $5$ urah, tretji v $12$ urah. Delavci
    skupaj začnejo zidati steno. Po eni uri drugi delavec odide, pridruži pa se četrti delavec.
    Skupaj s prvim in tretjim delavcem nato končajo steno v dveh urah. V kolikšnem času četrti delavec
    pozida steno?
    
    
    \begin{solution}[9cm]
        
        Zapis enačbe po besedilu $1$ točka, pravilno reševanje enačbe $1$ točka, zapis pravilne rešitve $1$ točka, zapis pravilnega odgovora $1$ točka.
    \end{solution}




        
\newpage
%%%% PKS, funkcija, premica

    % 1. vprašanje

    \question[4]
    Narišite dano množico točk v ravnini.
        $$[-1,2)\times \{-1,1,2\}$$

            \begin{solution}[9cm]
                ????
            \end{solution}

    
    
    % \newpage

    % 2. vprašanje
    \question[4]
    Izračunajte ploščino in orientacijo trikotnika $\triangle ABC$ z oglišči v točkah $A(4,-3)$, $B(5,4)$ in $C(-4,1)$.
    
    
    \begin{solution}[6cm]
        ???
    \end{solution}

        


\newpage

    % 5. vprašanje
    \question
    Dana je funkcija $$f(x)=\begin{cases}
    -2x-2, & x\leq 0 \\ -1, & 0<x<1 \\ 0.5x+1, & x\geq 2
    \end{cases}.$$ 
    
    \begin{parts}
        \part[6]
        Narišite dano funkcijo.

            \begin{solution}[11cm]
                ???
            \end{solution}

        \part[1]
        Določite definicijsko območje funkcije $f$.
        
            \begin{solution}[1cm]
                ???
            \end{solution}

        \part[3]
        Izračunajte ničle in začetno vrednost funkcije $f$.

            \begin{solution}[6cm]
                ???
            \end{solution}

        % \part[2]
        % Določite in utemeljite, ali je funkcija surjektivna. 
        
        %     \begin{solution}[3cm]
        %         ???
        %     \end{solution}

        \part[1]
        Določite interval, kjer je funkcija $f$ konstantna.
        
            \begin{solution}[0cm]
                ???
            \end{solution}

    \end{parts}
    
    






\end{questions}



\end{document}